\chapter{Configuración experimental}
\label{app:configuracion}

Este apéndice resume la configuración necesaria para reproducir los experimentos cerrados. Los valores proceden de los scripts y artefactos generados en el repositorio, no de un plan prospectivo.

\section{Entorno de software}

El proyecto declara Python 3.12 o posterior y utiliza \texttt{uv} como gestor de entorno. Las dependencias directas y opcionales relevantes se recogen en la tabla \ref{tab:dependencias_python_final}.

\begin{table}[H]
  \centering
  \caption{Dependencias principales declaradas por el proyecto.}
  \label{tab:dependencias_python_final}
  \small
  \begin{tabularx}{0.92\textwidth}{Xl}
    \toprule
    Paquete & Versión mínima \\
    \midrule
    Matplotlib & 3.9 \\
    NumPy & 2.0 \\
    OpenAI Python & 2.35.1 \\
    pandas & 2.2 \\
    Pillow & 10.4 \\
    Pydantic & 2.8 \\
    scikit-learn & 1.5 \\
    tqdm & 4.66 \\
    PyTorch, extra \texttt{cnn} & 2.3 \\
    Torchvision, extra \texttt{cnn} & 0.18 \\
    WFDB, extra \texttt{real-data} & 4.1 \\
    pytest, extra \texttt{dev} & 8.3 \\
    Ruff, extra \texttt{dev} & 0.6 \\
    \bottomrule
  \end{tabularx}
\end{table}

La memoria se compila con XeLaTeX a través del helper local de compilación. Los experimentos CNN se ejecutan con backend Matplotlib \texttt{Agg} para generar informes sin interfaz gráfica.

\section{Conjunto sintético QRS/ST}

\begin{table}[H]
  \centering
  \caption{Configuración del conjunto sintético final.}
  \label{tab:config_sintetico_final}
  \small
  \begin{tabularx}{0.94\textwidth}{p{4.4cm}X}
    \toprule
    Propiedad & Valor \\
    \midrule
    Fuente & \texttt{ecg\_few.simulator} \\
    Pacientes & 100 \\
    Imágenes & 300 \\
    Familias fuente & NORMAL, RBBB, ST\_ELEVATION, T\_WAVE\_INVERSION, COMBINED\_QRS\_ST \\
    Pacientes por familia & 20 \\
    Derivaciones renderizadas & V1, V2 y V3 \\
    Frecuencia de muestreo & 500 Hz \\
    Duración del beat & 900 ms \\
    Resolución de guardado & 130 ppp \\
    Positivos derivados & 20 \\
    Negativos o incompletos & 80 \\
    Presupuestos & \(K=\{2,4,8,16,32\}\) \\
    Semillas & 42, 123 y 2026 \\
    \bottomrule
  \end{tabularx}
\end{table}

La etiqueta positiva se obtiene cuando RBBB, elevación ST e inversión T están presentes simultáneamente. La agregación de paciente promedia las probabilidades de las derivaciones y después aplica el umbral seleccionado.

\section{Brugada HUCA}

\begin{table}[H]
  \centering
  \caption{Configuración del conjunto real procesado.}
  \label{tab:config_huca_final}
  \small
  \begin{tabularx}{0.94\textwidth}{p{4.4cm}X}
    \toprule
    Propiedad & Valor \\
    \midrule
    Fuente & Brugada HUCA 1.0.0 \\
    Pacientes brutos & 363 \\
    Pacientes excluidos & 46 por patrón basal limítrofe \\
    Pacientes incluidos & 317 \\
    Imágenes incluidas & 951 \\
    Casos confirmados & 116 \\
    Controles & 201 \\
    Derivaciones & V1, V2 y V3 \\
    Ventana & 300 ms antes de R y 600 ms después de R \\
    Frecuencia de origen & 100 Hz \\
    Resolución de guardado & 130 ppp \\
    Separación & Leave-one-out por paciente \\
    \bottomrule
  \end{tabularx}
\end{table}

HUCA no contiene anotaciones independientes para las tres etiquetas morfológicas del simulador. La tarea real utiliza la etiqueta clínica \texttt{clinical\_brugada}.

\section{Configuración CNN}

\begin{table}[H]
  \centering
  \caption{Hiperparámetros de la CNN final.}
  \label{tab:config_cnn_appendix}
  \small
  \begin{tabularx}{0.94\textwidth}{p{4.4cm}X}
    \toprule
    Propiedad & Valor \\
    \midrule
    Arquitectura & ResNet18 \\
    Pesos iniciales & \texttt{default} de Torchvision \\
    Entrada & RGB, 224 por 224 \\
    Salida & Un logit binario \\
    Pérdida & \texttt{BCEWithLogitsLoss} con \texttt{loss\_pos\_weight=auto} \\
    Optimizador & AdamW \\
    Tasa de aprendizaje & \(3\times10^{-4}\) \\
    Lote & 32 \\
    Épocas máximas & 20 \\
    Estrategia de umbral & \texttt{val\_derived\_balanced\_accuracy} \\
    Métrica principal & Exactitud equilibrada por paciente \\
    Grad CAM & 12 paneles en sintético y 6 en HUCA por ejecución agregada \\
    \bottomrule
  \end{tabularx}
\end{table}

\section{Configuración de adaptación de dominio}

\begin{table}[H]
  \centering
  \caption{Configuración de la rama de adaptación de dominio.}
  \label{tab:config_domain_appendix}
  \small
  \begin{tabularx}{0.94\textwidth}{p{4.4cm}X}
    \toprule
    Propiedad & Valor \\
    \midrule
    Origen etiquetado & \texttt{data/simulator\_qrs} \\
    Destino no etiquetado & \texttt{data/brugada\_huca} \\
    Métodos & SSL, CORAL, MMD y DANN \\
    Presupuestos & \(K=16\) y \(K=32\) \\
    Semillas & 42, 123 y 2026 \\
    Peso de adaptación & 0,1 \\
    Escalas MMD & 0,5, 1,0, 2,0 y 4,0 \\
    SSL & SimCLR durante 3 épocas, \texttt{ssl\_pretrain\_lr=0.0001} \\
    Política de destino & Transductiva compartida, sin etiquetas destino en la pérdida supervisada \\
    \bottomrule
  \end{tabularx}
\end{table}

\section{Configuración VLM/ICL}

La línea VLM/ICL forma parte del diseño comparativo al mismo nivel que la CNN. La configuración queda congelada para que las inferencias usen la misma base experimental, las mismas particiones por paciente y los mismos presupuestos \(K\):

\begin{itemize}
  \item modelos abiertos: \texttt{google/gemma-4-E4B-it} y \texttt{google/medgemma-4b-it}, configurables mediante \texttt{VLM\_MODELS};
  \item servidor local compatible con API de vLLM/OpenAI, configurado con \texttt{VLM\_API\_BASE};
  \item temperatura 0;
  \item salida JSON validada con las claves \texttt{RBBB}, \texttt{ST\_ELEVATION} y \texttt{T\_WAVE\_INVERSION};
  \item presupuestos \(K=\{0,2,4,8,16,32\}\);
  \item condiciones \texttt{zero\_shot}, \texttt{normal}, \texttt{balanced}, \texttt{permuted} y \texttt{no\_support\_images};
  \item salidas en \texttt{outputs/vlm\_loocv/} y \texttt{reports/loocv/vlm/}.
\end{itemize}

El capítulo~\ref{chap:vlm} ya incluye las tablas de salida, las condiciones experimentales y los criterios de comparación. Las celdas numéricas se rellenarán únicamente cuando existan predicciones por muestra, matrices de confusión y auditoría equivalente a la de CNN.
